% Part: first-order-logic
% Chapter: tableaux
% Section: proving-things

\documentclass[../../../include/open-logic-section]{subfiles}

\begin{document}

\iftag{FOL}
      {\olfileid{fol}{tab}{pro}}
      {\olfileid{pl}{tab}{pro}}

\olsection{\usetoken{P}{tableau}-এর উদাহরণ}

\begin{ex}
!!{sentence}~$(!A \land !B) \lif !A$-এর জন্য একটি বদ্ধ
!!{tableau} খুঁজি।

শুরুতে সংশ্লিষ্ট অনুমিতিটি !!{tableau}-এর একেবারে ওপরে লিখি।
\begin{oltableau}
  [\sFmla{\False}{(\formula{A} \land \formula{B}) \lif \formula{A}},
    just = \TAss]
\end{oltableau}

এখানে একটিই অনুমিতি, তাই বিধি প্রয়োগের উপযোগী !!{signed formula}-ও
একটি। (প্রতিটি !!{signed formula}-এ সর্বাধিক একটি বিধিই প্রয়োগ করা
যায়: সেটি হলো !!{sentence}-টির সংশ্লিষ্ট চিহ্ন ও !!{main operator}-এর
বিধি।) এখানে তাই $\TRule{\False}{\lif}$ প্রয়োগ করতেই হবে।
\begin{oltableau}
  [\sFmla{\False}{(\formula{A} \land \formula{B}) \lif \formula{A}},
    checked, just = \TAss
    [\sFmla{\True}{\formula{A} \land \formula{B}},
      just={\TRule{\False}{\lif}[1]}
      [\sFmla{\False}{\formula{A}}, just={\TRule{\False}{\lif}[1]}]
    ]
  ]
\end{oltableau}
কোন কোন !!{signed formula}s-এ সংশ্লিষ্ট বিধি প্রয়োগ করেছি তা মনে
রাখতে বাক্যের পাশে টিকচিহ্ন দিই। তবে কোনো বিধি সব খোলা শাখায়
প্রয়োগ করা হলেই \emph{কেবল} টিকচিহ্ন দেবে। কোনো !!a{signed formula}-এ
প্রতিটি খোলা শাখায় সংশ্লিষ্ট বিধি প্রয়োগ করা হয়ে গেলে আর ফিরে এসে
সেখানে বিধিটি প্রয়োগ করতে হবে না। এখানে একটি শাখাই আছে, তাই বিধিটি
মাত্র একবার প্রয়োগ করতে হয়েছে। (টিকচিহ্ন কেবল tableaux নির্মাণের
সুবিধা; tableaux-এর আনুষ্ঠানিক সংকেতবিন্যাসের অংশ নয়।)

এখন বিধি প্রয়োগের উপযোগী একটি নতুন !!{signed formula} আছে: পংক্তি~$2$-এর
$\sFmla{\True}{!A \land !B}$। $\TRule{\True}{\land}$ বিধি প্রয়োগে পাই:
\begin{oltableau}
  [\sFmla{\False}{(\formula{A} \land \formula{B}) \lif \formula{A}},
    checked, just = \TAss
    [\sFmla{\True}{\formula{A} \land \formula{B}},
      just={\TRule{\False}{\lif}[1]}, checked
      [\sFmla{\False}{\formula{A}}, just={\TRule{\False}{\lif}[1]}
        [\sFmla{\True}{\formula{A}}, just={\TRule{\True}{\land}[2]}
          [\sFmla{\True}{\formula{B}}, just={\TRule{\True}{\land}[2]}, close
          ]
        ]
      ]
    ]
  ]
\end{oltableau}
শাখাটিতে এখন পংক্তি~$4$-এর $\sFmla{\True}{!A}$ এবং পংক্তি~$3$-এর
$\sFmla{\False}{!A}$ উভয়ই আছে বলে শাখাটি বদ্ধ। এটিই একমাত্র শাখা,
তাই !!{tableau}-টিও বদ্ধ। ফলে $(!A \land !B) \lif !A$-এর জন্য
একটি বদ্ধ !!{tableau} পাওয়া গেল।
\end{ex}

\begin{ex}
এবার $(\lnot !A \lor !B) \lif (!A \lif !B)$-এর জন্য একটি বদ্ধ
!!{tableau} খুঁজি।

সংশ্লিষ্ট অনুমিতি দিয়ে শুরু করি:
\begin{oltableau}
  [\sFmla{\False}{(\lnot \formula{A} \lor \formula{B}) \lif
      (\formula{A} \lif \formula{B})}, just=\TAss]
\end{oltableau}
এই !!{tableau}-এর একমাত্র !!{signed formula}-টির
!!{main operator}~$\lif$ এবং চিহ্ন~$\False$; তাই এতে
$\TRule{\False}{\lif}$ বিধি প্রয়োগ করে পাই:
\begin{oltableau}
  [\sFmla{\False}{(\lnot \formula{A} \lor \formula{B})
      \lif (\formula{A} \lif \formula{B})}, just=\TAss, checked
    [\sFmla{\True}{\lnot \formula{A} \lor \formula{B}},
      just={\TRule{\False}{\lif}[1]}
      [\sFmla{\False}{(\formula{A} \lif \formula{B})},
        just={\TRule{\False}{\lif}[1]}
      ]
    ]
  ]
\end{oltableau}
এখন পংক্তি~$2$-এ $\TRule{\True}{\lor}$ প্রয়োগ করব, নাকি
পংক্তি~$3$-এ $\TRule{\False}{\lif}$ প্রয়োগ করব, তা বেছে নেওয়া যায়।
প্রতিটি !!{signed formula}-এ তার বিধি প্রত্যেক শাখায় প্রয়োগ করা হলে
ক্রমটি আসলে গুরুত্বপূর্ণ নয়। তাই প্রথমটিই নিই। $\TRule{\True}{\lor}$
বিধি !!{tableau}-কে শাখায় ভাগ করতে দেয়; বিধিটির দুটি সিদ্ধান্ত দুই
নতুন শাখায় যোগ হওয়া নতুন !!{signed formula}s হবে। ফলে পাই:
\begin{oltableau}
  [\sFmla{\False}{(\lnot \formula{A} \lor \formula{B}) \lif
      (\formula{A} \lif \formula{B})}, just=\TAss, checked
    [\sFmla{\True}{\lnot \formula{A} \lor \formula{B}},
      just={\TRule{\False}{\lif}[1]}, checked
      [\sFmla{\False}{(\formula{A} \lif \formula{B})},
        just={\TRule{\False}{\lif}[1]}
        [\sFmla{\True}{\lnot \formula{A}}, just={\TRule{\True}{\lor}[2]}]
        [\sFmla{\True}{\formula{B}}, just={\TRule{\True}{\lor}[2]}]
      ]
    ]
  ]
\end{oltableau}
পংক্তি~$3$-এ এখনও $\TRule{\False}{\lif}$ বিধি প্রয়োগ করা হয়নি;
এবার তা করি। সময় বাঁচাতে দুই শাখাতেই একসঙ্গে প্রয়োগ করি।
মনে রেখো, কোনো !!a{signed formula}-এর পাশে তখনই টিকচিহ্ন দিই, যখন
প্রতিটি খোলা শাখায় সংশ্লিষ্ট বিধিটি প্রয়োগ করা হয়েছে। তাই বিধিটি
যে !!{signed formula}-এ প্রযোজ্য, সেটি-সম্বলিত প্রত্যেক শাখার শেষে
বিধিটি প্রয়োগ করা ভালো। তাহলে নানা শাখার আরও নিচে গিয়ে সেই
!!{signed formula}-টিতে আবার ফিরতে হবে না।
\begin{oltableau}
  [\sFmla{\False}{(\lnot \formula{A} \lor \formula{B}) \lif
      (\formula{A} \lif \formula{B})}, just=\TAss, checked
    [\sFmla{\True}{\lnot \formula{A} \lor \formula{B}},
      just={\TRule{\False}{\lif}[1]}, checked
      [\sFmla{\False}{(\formula{A} \lif \formula{B})},
        just={\TRule{\False}{\lif}[1]}, checked
        [\sFmla{\True}{\lnot \formula{A}}, just={\TRule{\True}{\lor}[2]}
          [\sFmla{\True}{\formula{A}}, just={\TRule{\False}{\lif}[3]}
            [\sFmla{\False}{\formula{B}}, just={\TRule{\False}{\lif}[3]}]
          ]
        ]
        [\sFmla{\True}{\formula{B}}, just={\TRule{\True}{\lor}[2]}
          [\sFmla{\True}{\formula{A}}, just={\TRule{\False}{\lif}[3]}
            [\sFmla{\False}{\formula{B}}, just={\TRule{\False}{\lif}[3]}, close]
          ]
        ]
      ]
    ]
  ]
\end{oltableau}
ডান শাখাটি এখন বদ্ধ। বাঁ শাখায় পংক্তি~$4$-এ এখনও
$\TRule{\True}{\lnot}$ বিধি প্রয়োগ করা যায়। এতে
$\sFmla{\False}{!A}$ পাওয়া যায় এবং বাঁ শাখাটিও বদ্ধ হয়:
\begin{oltableau}
  [\sFmla{\False}{(\lnot \formula{A} \lor \formula{B}) \lif
      (\formula{A} \lif \formula{B})}, just=\TAss, checked
    [\sFmla{\True}{\lnot \formula{A} \lor \formula{B}},
      just={\TRule{\False}{\lif}[1]}, checked
      [\sFmla{\False}{(\formula{A} \lif \formula{B})},
        just={\TRule{\False}{\lif}[1]}, checked
        [\sFmla{\True}{\lnot \formula{A}}, just={\TRule{\True}{\lor}[2]}
          [\sFmla{\True}{\formula{A}}, just={\TRule{\False}{\lif}[3]}
            [\sFmla{\False}{\formula{B}}, just={\TRule{\False}{\lif}[3]}
              [\sFmla{\False}{\formula{A}},
                just={\TRule{\True}{\lnot}[4]}, close]
            ]
          ]
        ]
        [\sFmla{\True}{\formula{B}}, just={\TRule{\True}{\lor}[2]}
          [\sFmla{\True}{\formula{A}}, just={\TRule{\False}{\lif}[3]}
            [\sFmla{\False}{\formula{B}},
              just={\TRule{\False}{\lif}[3]}, close]
          ]
        ]
      ]
    ]
  ]
\end{oltableau}
\end{ex}

\begin{ex}
যেকোনো সংখ্যক !!{signed formula}s-কে অনুমিতি ধরে !!{tableau}s দেওয়া
যায়। অনেক সময় একাধিক শাখা-বিভাজক বিধি প্রয়োগ করাও দরকার; সাধারণভাবে
!!a{tableau}-তে যেকোনো সংখ্যক শাখা থাকতে পারে। যেমন,
$\{\sFmla{\True}{!A \lor (!B \land !C)}, \sFmla{\False}{(!A \lor !B) \land (!A
  \lor !C)}\}$-এর জন্য !!a{tableau} বিবেচনা করো। প্রথম অনুমিতিতে
$\TRule{\True}{\lor}$ প্রয়োগ করে শুরু করি:
\begin{oltableau}
  [\sFmla{\True}{\formula{A} \lor (\formula{B} \land \formula{C})},
    just=\TAss, checked
    [\sFmla{\False}{(\formula{A} \lor \formula{B}) \land
        (\formula{A} \lor \formula{C})}, just=\TAss
      [\sFmla{\True}{\formula{A}}, just={\TRule{\True}{\lor}[1]}]
      [\sFmla{\True}{\formula{B} \land \formula{C}},
        just={\TRule{\True}{\lor}[1]}]
    ]
  ]
\end{oltableau}
এবার পংক্তি~$2$-এ $\TRule{\False}{\land}$ বিধি প্রয়োগ করা যায়।
দুই শাখায় একসঙ্গে তা করি, তাই পংক্তি~$2$-এ টিকচিহ্নও দেওয়া যায়:
\begin{oltableau}
  [\sFmla{\True}{\formula{A} \lor (\formula{B} \land \formula{C})},
    just=\TAss, checked
    [\sFmla{\False}{(\formula{A} \lor \formula{B}) \land
        (\formula{A} \lor \formula{C})}, just=\TAss, checked
      [\sFmla{\True}{\formula{A}}, just={\TRule{\True}{\lor}[1]}
        [\sFmla{\False}{\formula{A} \lor \formula{B}},
          just={\TRule{\False}{\land}[2]}]
        [\sFmla{\False}{\formula{A} \lor \formula{C}},
          just={\TRule{\False}{\land}[2]}]
      ]
      [\sFmla{\True}{\formula{B} \land \formula{C}},
        just={\TRule{\True}{\lor}[1]}
        [\sFmla{\False}{\formula{A} \lor \formula{B}},
          just={\TRule{\False}{\land}[2]}]
        [\sFmla{\False}{\formula{A} \lor \formula{C}},
          just={\TRule{\False}{\land}[2]}]
      ]
    ]
  ]
\end{oltableau}
এখন $!A \lor !B$-সম্বলিত সব শাখায়
$\TRule{\False}{\lor}$ প্রয়োগ করি:
\begin{oltableau}
  [\sFmla{\True}{\formula{A} \lor (\formula{B} \land \formula{C})},
    just=\TAss, checked
    [\sFmla{\False}{(\formula{A} \lor \formula{B}) \land
        (\formula{A} \lor \formula{C})}, just=\TAss, checked
      [\sFmla{\True}{\formula{A}}, just={\TRule{\True}{\lor}[1]}
        [\sFmla{\False}{\formula{A} \lor \formula{B}},
          just={\TRule{\False}{\land}[2]}, checked
          [\sFmla{\False}{\formula{A}}, just={\TRule{\False}{\lor}[4]}
            [\sFmla{\False}{\formula{B}},
              just={\TRule{\False}{\lor}[4]}, close]
          ]
        ]
        [\sFmla{\False}{\formula{A} \lor \formula{C}},
          just={\TRule{\False}{\land}[2]}]
      ]
      [\sFmla{\True}{\formula{B} \land \formula{C}},
        just={\TRule{\True}{\lor}[1]}
        [\sFmla{\False}{\formula{A} \lor \formula{B}},
          just={\TRule{\False}{\land}[2]}, checked
          [\sFmla{\False}{\formula{A}}, just={\TRule{\False}{\lor}[4]}
            [\sFmla{\False}{\formula{B}},
              just={\TRule{\False}{\lor}[4]}]
          ]
        ]
        [\sFmla{\False}{\formula{A} \lor \formula{C}},
          just={\TRule{\False}{\land}[2]}]
      ]
    ]
  ]
\end{oltableau}
সবচেয়ে বাঁয়ের শাখাটি এখন বদ্ধ। এবার $!A \lor !C$-তে
$\TRule{\False}{\lor}$ প্রয়োগ করি:
\begin{oltableau}
  [\sFmla{\True}{\formula{A} \lor (\formula{B} \land \formula{C})},
    just=\TAss, checked
    [\sFmla{\False}{(\formula{A} \lor \formula{B}) \land
        (\formula{A} \lor \formula{C})}, just=\TAss, checked
      [\sFmla{\True}{\formula{A}}, just={\TRule{\True}{\lor}[1]}
        [\sFmla{\False}{\formula{A} \lor \formula{B}},
          just={\TRule{\False}{\land}[2]}, checked
          [\sFmla{\False}{\formula{A}},
            just={\TRule{\False}{\lor}[4]}
            [\sFmla{\False}{\formula{B}},
              just={\TRule{\False}{\lor}[4]}, close]
          ]
        ]
        [\sFmla{\False}{\formula{A} \lor \formula{C}},
          just={\TRule{\False}{\land}[2]}, checked
          [\sFmla{\False}{\formula{A}},
            just={\TRule{\False}{\lor}[4]}, move by=2
            [\sFmla{\False}{\formula{C}},
              just={\TRule{\False}{\lor}[4]}, close]
          ]
        ]
      ]
      [\sFmla{\True}{\formula{B} \land \formula{C}},
        just={\TRule{\True}{\lor}[1]}
        [\sFmla{\False}{\formula{A} \lor \formula{B}},
          just={\TRule{\False}{\land}[2]}, checked
          [\sFmla{\False}{\formula{A}},
            just={\TRule{\False}{\lor}[4]}
            [\sFmla{\False}{\formula{B}},
              just={\TRule{\False}{\lor}[4]}
            ]
          ]
        ]
        [\sFmla{\False}{\formula{A} \lor \formula{C}},
          just={\TRule{\False}{\land}[2]}, checked
          [\sFmla{\False}{\formula{A}},
            just={\TRule{\False}{\lor}[4]}, move by=2
            [\sFmla{\False}{\formula{C}},
              just={\TRule{\False}{\lor}[4]}]
          ]
        ]
      ]
    ]
  ]
\end{oltableau}
স্পষ্টতার জন্য $\TRule{\False}{\lor}$-এর দ্বিতীয় প্রয়োগের ফলটি
নিচে সরানো হয়েছে। এখানে তার প্রয়োজন ছিল না, কারণ সমর্থন একই থাকত।

দুটি শাখা এখনও খোলা এবং পংক্তি~$3$-এর
$\sFmla{\True}{!B \land !C}$-তে টিকচিহ্ন নেই। এতে
$\TRule{\True}{\land}$ প্রয়োগ করে একটি বদ্ধ !!{tableau} পাই:
\begin{oltableau}
  [\sFmla{\True}{\formula{A} \lor (\formula{B} \land \formula{C})},
    just=\TAss, checked
    [\sFmla{\False}{(\formula{A} \lor \formula{B}) \land
        (\formula{A} \lor \formula{C})}, just=\TAss, checked
      [\sFmla{\True}{\formula{A}}, just={\TRule{\True}{\lor}[1]}
        [\sFmla{\False}{\formula{A} \lor \formula{B}},
          just={\TRule{\False}{\land}[2]}, checked
          [\sFmla{\False}{\formula{A}},
            just={\TRule{\False}{\lor}[4]}
            [\sFmla{\False}{\formula{B}},
              just={\TRule{\False}{\lor}[4]}, close]
          ]
        ]
        [\sFmla{\False}{\formula{A} \lor \formula{C}},
          just={\TRule{\False}{\land}[2]}, checked
          [\sFmla{\False}{\formula{A}},
            just={\TRule{\False}{\lor}[4]}
            [\sFmla{\False}{\formula{C}},
              just={\TRule{\False}{\lor}[4]}, close]
          ]
        ]
      ]
      [\sFmla{\True}{\formula{B} \land \formula{C}},
        just={\TRule{\True}{\lor}[1]}, checked
        [\sFmla{\False}{\formula{A} \lor \formula{B}},
          just={\TRule{\False}{\land}[2]}, checked
          [\sFmla{\False}{\formula{A}},
            just={\TRule{\False}{\lor}[4]}
            [\sFmla{\False}{\formula{B}},
              just={\TRule{\False}{\lor}[4]}
              [\sFmla{\True}{\formula{B}},
                just={\TRule{\True}{\land}[3]}
                [\sFmla{\True}{\formula{C}},
                  just={\TRule{\True}{\land}[3]},close
                ]
              ]
            ]
          ]
        ]
        [\sFmla{\False}{\formula{A} \lor \formula{C}},
          just={\TRule{\False}{\land}[2]}, checked
          [\sFmla{\False}{\formula{A}},
            just={\TRule{\False}{\lor}[4]}
            [\sFmla{\False}{\formula{C}},
              just={\TRule{\False}{\lor}[4]}
              [\sFmla{\True}{\formula{B}},
                just={\TRule{\True}{\land}[3]}
                [\sFmla{\True}{\formula{C}},
                  just={\TRule{\True}{\land}[3]},close
                ]
              ]
            ]
          ]
        ]
      ]
    ]
  ]
\end{oltableau}

তুলনার জন্য একই অনুমিতি-সমষ্টির একটি বদ্ধ !!{tableau} নিচে দেওয়া
হলো, যেখানে বিধিগুলি অন্য ক্রমে প্রয়োগ করা হয়েছে:
\begin{oltableau}
  [\sFmla{\True}{\formula{A} \lor (\formula{B} \land \formula{C})},
    just=\TAss, checked
    [\sFmla{\False}{(\formula{A} \lor \formula{B}) \land
        (\formula{A} \lor \formula{C})}, just=\TAss, checked
      [\sFmla{\False}{\formula{A} \lor \formula{B}},
        just={\TRule{\False}{\land}[2]}, checked
        [\sFmla{\False}{\formula{A}},
          just={\TRule{\False}{\lor}[3]}
          [\sFmla{\False}{\formula{B}},
            just={\TRule{\False}{\lor}[3]}
            [\sFmla{\True}{\formula{A}},
              just={\TRule{\True}{\lor}[1]},close
            ]
            [\sFmla{\True}{\formula{B} \land \formula{C}},
              just={\TRule{\True}{\lor}[1]}, checked
              [\sFmla{\True}{\formula{B}},
                just={\TRule{\True}{\land}[6]}
                [\sFmla{\True}{\formula{C}},
                  just={\TRule{\True}{\land}[6]},close
                ]
              ]
            ]
          ]
        ]
      ]
      [\sFmla{\False}{\formula{A} \lor \formula{C}},
        just={\TRule{\False}{\land}[2]}, checked
        [\sFmla{\False}{\formula{A}},
          just={\TRule{\False}{\lor}[3]}
          [\sFmla{\False}{\formula{C}},
            just={\TRule{\False}{\lor}[3]}
            [\sFmla{\True}{\formula{A}},
              just={\TRule{\True}{\lor}[1]},close
            ]
            [\sFmla{\True}{\formula{B} \land \formula{C}},
              just={\TRule{\True}{\lor}[1]}, checked
              [\sFmla{\True}{\formula{B}},
                just={\TRule{\True}{\land}[6]}
                [\sFmla{\True}{\formula{C}},
                  just={\TRule{\True}{\land}[6]},close
                ]
              ]
            ]
          ]
        ]
      ]
    ]
  ]
\end{oltableau}
\end{ex}

\begin{prob}
নিচেরগুলির বদ্ধ !!{tableau}s দাও:
\begin{enumerate}
\item $\sFmla{\True}{!A \land (!B \land !C)}, \sFmla{\False}{(!A \land !B) \land !C}$।
\item $\sFmla{\True}{!A \lor (!B \lor !C)}, \sFmla{\False}{(!A \lor !B) \lor !C}$।
\item $\sFmla{\True}{!A \lif (!B \lif !C)}, \sFmla{\False}{!B \lif (!A \lif !C)}$।
\item $\sFmla{\True}{!A}, \sFmla{\False}{\lnot\lnot !A}$।
\end{enumerate}
\end{prob}

\begin{prob}
নিচেরগুলির বদ্ধ !!{tableau}s দাও:
\begin{enumerate}
\item $\sFmla{\True}{(!A \lor !B) \lif !C}, \sFmla{\False}{!A \lif !C}$।
\item $\sFmla{\True}{(!A \lif !C) \land (!B \lif !C)}, \sFmla{\False}{(!A \lor !B) \lif !C}$।
\item $\sFmla{\False}{\lnot(!A \land \lnot !A)}$।
\item $\sFmla{\True}{!B \lif !A}, \sFmla{\False}{\lnot !A \lif \lnot !B}$।
\item $\sFmla{\False}{(!A \lif \lnot !A) \lif \lnot !A}$।
\item $\sFmla{\False}{\lnot(!A \lif !B) \lif \lnot !B}$।
\item $\sFmla{\True}{!A \lif !C}, \sFmla{\False}{\lnot (!A \land \lnot !C)}$।
\item $\sFmla{\True}{!A \land \lnot !C}, \sFmla{\False}{\lnot (!A \lif !C)}$।
\item $\sFmla{\True}{!A \lor !B}, \sFmla{\True}{\lnot !B}, \sFmla{\False}{!A}$।
\item $\sFmla{\True}{\lnot !A \lor \lnot !B}, \sFmla{\False}{\lnot(!A \land !B)}$।
\item $\sFmla{\False}{(\lnot !A \land \lnot !B) \lif\lnot(!A \lor !B)}$।
\item $\sFmla{\False}{\lnot(!A \lor !B) \lif (\lnot !A \land \lnot !B)}$।
\end{enumerate}
% BN-SRC-046 TARGET-CORRECTION-BEGIN
\par\smallskip\noindent\textnormal{\small\textbf{উৎস-সংশোধন (BN-SRC-046):} হিমায়িত ইংরেজি নবম অনুশীলনে $!A \lor !B$ ও $\lnot !B$-কে ভুল করে একই \texttt{\string\sFmla}-র সূত্র-ঘরে কমা দিয়ে রাখা হয়েছে। অর্থপূর্ণ তিনটি চিহ্নিত সূত্র হিসেবে বাংলা লক্ষ্যে $\sFmla{\True}{!A \lor !B}$, $\sFmla{\True}{\lnot !B}$ ও $\sFmla{\False}{!A}$ আলাদা করা হয়েছে।} % BN-SRC-046 TARGET-CORRECTION-END
\end{prob}

\begin{prob}
নিচেরগুলির বদ্ধ !!{tableau}s দাও:
\begin{enumerate}
\item $\sFmla{\True}{\lnot(!A \lif !B)}, \sFmla{\False}{!A}$।
\item $\sFmla{\True}{\lnot(!A \land !B)}, \sFmla{\False}{\lnot !A \lor \lnot !B}$।
\item $\sFmla{\True}{!A \lif !B}, \sFmla{\False}{\lnot !A \lor !B}$।
\item $\sFmla{\False}{\lnot \lnot !A \lif !A}$।
\item $\sFmla{\True}{!A \lif !B}, \sFmla{\True}{\lnot !A \lif !B}, \sFmla{\False}{!B}$।
\item $\sFmla{\True}{(!A \land !B) \lif !C}, \sFmla{\False}{(!A \lif !C) \lor (!B \lif !C)}$।
\item $\sFmla{\True}{(!A \lif !B) \lif !A}, \sFmla{\False}{!A}$।
\item $\sFmla{\False}{(!A \lif !B) \lor (!B \lif !C)}$।
\end{enumerate}
\end{prob}

\end{document}
